%
% Niniejszy plik stanowi przykład formatowania pracy magisterskiej na
% Wydziale MIM UW.  Szkielet użytych poleceń można wykorzystywać do
% woli, np. formatujac wlasna prace.
%
% Zawartosc merytoryczna stanowi oryginalnosiagniecie
% naukowosciowe Marcina Wolinskiego.  Wszelkie prawa zastrzeżone.
%
% Copyright (c) 2001 by Marcin Woliński <M.Wolinski@gust.org.pl>
% Poprawki spowodowane zmianami przepisów - Marcin Szczuka, 1.10.2004
% Poprawki spowodowane zmianami przepisow i ujednolicenie 
% - Seweryn Karłowicz, 05.05.2006
% Dodanie wielu autorów i tłumaczenia na angielski - Kuba Pochrybniak, 29.11.2016

% dodaj opcję [licencjacka] dla pracy licencjackiej
% dodaj opcję [en] dla wersji angielskiej (mogą być obie: [licencjacka,en])
% \documentclass[magisterska]{pracamgr}
\documentclass{pracamgr}

% Dodatkowe importy
\usepackage[backend=biber, style=numeric, sortcites=true]{biblatex}
\usepackage{hyperref}
\usepackage{csquotes}
\addbibresource{Latex/bibliography.bib}


% Dane magistranta:
\autor{Michał Koziński}{464410}


% Dane magistrantów:
%\autor{Autor Zerowy}{342007}
%\autori{Autor Pierwszy}{342013}
%\autorii{Drugi Autor-Z-Rzędu}{231023}
%\autoriii{Trzeci z Autorów}{777321}
%\autoriv{Autor nr Cztery}{432145}
%\autorv{Autor nr Pięć}{342011}

\title{Holistyczna analiza obszarowa dostępności transportu zbiorowego w aglomeracjach miejskich.}


\tytulang{
%An implementation of a difference blabalizer based on the theory of $\sigma$ -- $\rho$ phetors
}

%kierunek: 
% - matematyka, informacyka, ...
% - Mathematics, Computer Science, ...
\kierunek{Inżynieria Obliczeniowa}

% informatyka - nie okreslamy zakresu (opcja zakomentowana)
% matematyka - zakres moze pozostac nieokreslony,
% a jesli ma byc okreslony dla pracy mgr,
% to przyjmuje jedna z wartosci:
% {metod matematycznych w finansach}
% {metod matematycznych w ubezpieczeniach}
% {matematyki stosowanej}
% {nauczania matematyki}
% Dla pracy licencjackiej mamy natomiast
% mozliwosc wpisania takiej wartosci zakresu:
% {Jednoczesnych Studiow Ekonomiczno--Matematycznych}

% \zakres{Tu wpisac, jesli trzeba, jedna z opcji podanych wyzej}

% Praca wykonana pod kierunkiem:
% (podać tytuł/stopień imię i nazwisko opiekuna
% Instytut
% ew. Wydział ew. Uczelnia (jeżeli nie MIM UW))
\opiekun{dr. hab. Dominika Batorskiego\\
  ICM, UW\\
  }

% miesiąc i~rok:
\date{Maj 2026}

%Podać dziedzinę wg klasyfikacji Socrates-Erasmus:
% \dziedzina{ 
% %11.0 Matematyka, Informatyka:\\ 
% 11.1 Matematyka\\ 
% %11.3 Informatyka\\ 
% }
 \dziedzina{?}

%Klasyfikacja tematyczna wedlug AMS (matematyka) lub ACM (informatyka)
% \klasyfikacja{D. Software\\
%   D.127. Blabalgorithms\\
%   D.127.6. Numerical blabalysis}
\klasyfikacja{?}

% Słowa kluczowe:
% \keywords{blabaliza różnicowa, fetory $\sigma$-$\rho$, fooizm,
%   blarbarucja, blaba, fetoryka, baleronik}
\keywords{?}

% Tu jest dobre miejsce na Twoje własne makra i~środowiska:
\newtheorem{defi}{Definicja}[section]

% koniec definicji

\begin{document}

\maketitle

%tu idzie streszczenie na strone poczatkowa
\begin{abstract}
  % W~pracy przedstawiono prototypową implementację blabalizatora
  % różnicowego bazującą na teorii fetorów $\sigma$-$\rho$ profesora
  % Fifaka.  Wykorzystanie teorii Fifaka daje wreszcie możliwość
  % efektywnego wykonania blabalizy numerycznej.  Fakt ten stanowi
  % przełom technologiczny, którego konsekwencje trudno z~góry
  % przewidzieć.
\end{abstract}

\tableofcontents
%\listoffigures
%\listoftables

\chapter*{Wstęp}
\addcontentsline{toc}{chapter}{Wstęp}

\chapter{Przegląd Literatury i Podstawy Metodologiczne}\label{chapter:literatura}

\section{Teoretyczne ramy dostępności: cztery komponenty dostępności}\label{section:komponenty}

W paradygmacie zaawansowanej analizy przestrzennej i~inżynierii transportu, dostępność (ang.~accessibility) nie jest wyłącznie pochodną odległości fizycznej, lecz stanowi fundamentalny \textit{produkt} zintegrowanego systemu transportowego, decydujący o przewadze lokalizacyjnej danego obszaru~\cite{turoń_sierpiński_2019}. W~ujęciu formalnym, dostępność definiuje się jako stopień, w~jakim połączone systemy zagospodarowania przestrzennego i~transportu masowego umożliwiają określonym grupom jednostek dotarcie do pożądanych celów aktywności (takich jak miejsca pracy, edukacji czy szpitale) przy użyciu dostępnych środków przemieszczania się w~określonym czasie~\cite{capetown}. Aby dogłębnie diagnozować złożone interakcje sieciowe, ewaluacja dostępności opiera się na rygorystycznych ramach metodologicznych, które dekomponują to zjawisko na cztery współzależne komponenty: transportowy, przestrzenny, czasowy oraz indywidualny~\cite{gravity2, literature_review}.

\subsection{Komponent transportowy (podażowy i infrastrukturalny)}\label{subsection:komponent_transportowy}
Komponent ten stanowi rdzeń podażowy całego systemu przemieszczania się, definiując fizyczne oraz operacyjne możliwości pokonywania odległości pomiędzy strefami źródłowymi a docelowymi. W najszerszym ujęciu odzwierciedla on wielkość oporu przestrzeni, z jakim musi zmierzyć się użytkownik sieci transportowej. Na wymiar ten składają się dwa nierozerwalnie połączone aspekty: twarda infrastruktura (układ dróg, torowisk, lokalizacja przystanków i węzłów przesiadkowych) oraz nałożona na nią miękka warstwa operacyjna (przebieg linii komunikacyjnych, parametry taboru, prędkości eksploatacyjne czy pojemność systemu)~\cite{mrt, l-space2, c-space}.

Z perspektywy analitycznej, komponent ten pełni rolę głównego pośrednika determinującego wydajność obsługi obszaru miejskiego. Jego rygorystyczna ocena prowadzona jest dwutorowo. Z jednej strony, strukturę i układ relacji wewnątrz systemu bada się przy użyciu aparatury matematycznej wywodzącej się z teorii grafów, co pozwala na identyfikację wąskich gardeł i ewaluację spójności powiązań. Z drugiej strony, z punktu widzenia pasażera, organizacja transportu bezpośrednio przekłada się na całkowity wysiłek i czas niezbędny do pokonania trasy. Wymusza to syntetyczną kwantyfikację tzw. impedancji, czyli uogólnionego kosztu podróży, który uwzględnia wszystkie obciążenia związane z przemieszczaniem się. W konsekwencji komponent transportowy nie tylko warunkuje fizyczną możliwość dotarcia do celu, ale poprzez uciążliwość samej podróży bezpośrednio kształtuje postrzeganą dostępność i atrakcyjność danego obszaru~\cite{capetown, LUO2019102505}.

\subsection{Komponent przestrzenny (land-use)}\label{subsection:komponent_przestrzenny}
Z drugiej strony równania dostępności znajduje się wymiar przestrzenny. Odnosi się on bezpośrednio do rozmieszczenia, wielkości oraz zróżnicowania rynków celu w tkance miejskiej (takich jak miejsca pracy, handlu, ochrony zdrowia czy edukacji)~\cite{elavating, capetown, Koźlak_Aleksandra_Znaczenie_2010}. W modelach transportowych obszar ten reprezentuje miarę atrakcyjności poszczególnych stref, stanowiąc docelowy punkt podróży użytkowników sieci. Współczesne ujęcie tego zagadnienia odchodzi od traktowania celów wyłącznie jako izolowanych punktów na mapie, zwracając uwagę na ograniczoną pojemność tych miejsc oraz zjawisko wzajemnej konkurencji o zasoby (np. wakaty na rynku pracy). Precyzyjne bilansowanie popytu z faktycznymi możliwościami rynków docelowych stanowi wyzwanie, które w zaawansowanej analityce rozwiązuje się poprzez implementację odpowiednich modeli grawitacyjnych, co zostało szerzej omówione w dalszej części pracy~\cite{gravity2, mrt}.

\subsection{Komponent czasowy}\label{subsection:komponent_czasowy}
Oprócz fizycznego pokonywania odległości, krytycznym aspektem ewaluacji systemów transportowych jest ich zmienność. Dostępność, ze szczególnym uwzględnieniem komunikacji zbiorowej, rzadko bywa wartością stałą – wykazuje za to znaczną fluktuację w zależności od pory dnia czy tygodnia. Wymiar czasowy spina ze sobą ograniczenia nakładane z jednej strony przez harmonogramy funkcjonowania samych celów podróży (np. godziny otwarcia placówek czy ramy czasu pracy), a z drugiej – przez rozkłady jazdy, które bezpośrednio determinują częstotliwość i niezawodność usług przewozowych~\cite{gravity2, capetown, review, LUO2019102505}. Zrozumienie tej dynamiki uświadamia, że obszary o rzekomo doskonałej łączności w godzinach szczytu mogą cierpieć na drastyczne wykluczenie wieczorami. Rygorystyczne ujęcie tych wahań wymaga włączenia do analizy czasu oczekiwania na przystanku, co zostanie rozwinięte w rozdziale poświęconym modelowaniu całego łańcucha podróży~\cite{capetown, LUO2019102505, elavating}.

\subsection{Komponent indywidualny (socjo-ekonomiczny i fizyczny)}\label{subsection:komponent_indywidualny}
Całościowy obraz analizowanego zjawiska domyka perspektywa samego pasażera. Należy bowiem podkreślić, że obiektywnie ta sama sieć może oferować diametralnie różny poziom możliwości przemieszczania się w zależności od cech i ograniczeń konkretnego użytkownika~\cite{gravity2, capetown, mrt}. Aspekt ten nakazuje ewaluację systemu przez pryzmat barier napotykanych przez poszczególne grupy demograficzne. Obejmuje to zarówno bariery fizyczne i architektoniczne, utrudniające podróżowanie osobom z niepełnosprawnościami, jak i uwarunkowania socjoekonomiczne, wśród których kluczowa staje się przystępność cenowa usług~\cite{capetown, elavating}. Z perspektywy systemowej, nawet najlepiej zaprojektowana infrastruktura staje się bezużyteczna, jeśli koszt przejazdu przewyższa budżet mieszkańców. Zjawisko to i wynikająca z niego konieczność segmentacji podróżnych stanowią fundament zagadnień sprawiedliwości transportowej, poruszanych w kolejnych sekcjach~\cite{capetown}.


\section{Topologiczna analiza sieci}\label{section:analiza_topologiczna}

Aby dokonać operacjonalizacji opisanego wcześniej komponentu transportowego, konieczne jest dobranie odpowiedniego aparatu matematycznego. Jak zasygnalizowano, rygorystyczna ocena wydajności transportu zbiorowego wymaga mapowania rzeczywistych układów przy użyciu teorii grafów. Nowoczesne podejście do modelowania złożonych systemów transportowych rozwiązuje ten problem, dekomponując ich strukturę na odrębne, lecz ściśle powiązane warstwy topologiczne: L-Space oraz P-Space. Zrozumienie tego dualizmu (oddzielenia fizycznej infrastruktury od logiki usług) stanowi kluczowy warunek w inżynierii transportu do identyfikowania wąskich gardeł i błędów operacyjnych w zarządzaniu ruchem~\cite{l-space2}.

\subsection{Przestrzeń infrastruktury (L-Space)}\label{subsection:l-space}
Topologia L-space stanowi podstawową reprezentację sieciową, w której węzły odpowiadają fizycznym przystankom lub stacjom, natomiast łączące je krawędzie definiuje się wyłącznie między przystankami następującymi bezpośrednio po sobie na trasie danej linii komunikacyjnej. Model ten jest w istocie wiernym odwzorowaniem fizycznej infrastruktury miejskiej – układu ulic, torowisk i szlaków osadzonych w topografii terenu. Z perspektywy inżynierii transportu badanie przestrzeni L-space okazuje się niezastąpione w diagnozowaniu problemów o charakterze podażowym. Pozwala ono na precyzyjną identyfikację wąskich gardeł (ang. \textit{bottlenecks}), estymację rzeczywistych odległości przestrzennych oddzielających węzły oraz ocenę fizycznej odporności (ang. \textit{vulnerability}) całego układu na awarie i zakłócenia w ruchu~\cite{c-space, l-space2}.

\subsection{Przestrzeń usług (P-space)}\label{subsection:p-space}
Z kolei topologia P-space oferuje całkowicie odmienną, zorientowaną na pasażera perspektywę analityczną. Chociaż zbiór węzłów nadal opiera się na fizycznych przystankach, to krawędzie łączą w niej wybrane strefy wtedy, gdy można między nimi podróżować bezpośrednio – przy użyciu pojedynczej linii komunikacyjnej, bez konieczności przesiadania się. Taka konstrukcja prowadzi do kluczowej własności matematycznej: w modelu P-space wszystkie przystanki obsługiwane przez konkretną linię tworzą pełny podgraf (klikę). Reprezentacja ta doskonale odzwierciedla operacyjne postrzeganie układu przez użytkownika. W tym ujęciu odległość topologiczna, czyli najkrótsza ścieżka wyznaczona w grafie, nie określa dystansu w kilometrach, lecz bezpośrednio kwantyfikuje minimalną liczbę linii niezbędnych do osiągnięcia celu podróży~\cite{c-space, l-space2}.

\subsection{Istotność separacji infrastruktury i usług}\label{subsection:serparacja}
Znaczenie analitycznego dualizmu L-space/P-space polega na ścisłym oddzieleniu infrastruktury (tego, co fizycznie istnieje w przestrzeni) od nałożonego na nią schematu operacyjnego (tego, jak owa przestrzeń jest wykorzystywana poprzez siatkę połączeń). Dopiero zestawienie ze sobą wskaźników sieciowych z obu tych topologii pozwala ujawnić utajone dysfunkcje układu.

Przykładowo, jeżeli dany system transportowy charakteryzuje się wysoką spójnością w topologii L-space, ale jednocześnie wykazuje duże średnie długości ścieżek w przestrzeni P-space, implikuje to błędny model zarządzania operacyjnego. Oznacza to bowiem, że pomimo doskonałej fizycznej siatki ulic czy torów, organizator narzucił na nią wysoce pofragmentowane usługi (np. bardzo krótkie trasy linii), zmuszając pasażerów do ciągłych przesiadek. Taka sytuacja drastycznie zwiększa uciążliwość przemieszczania się – każdy wymuszony transfer stanowi istotną barierę dla pasażera, która winduje całkowity wysiłek wkładany w podróż (co zostanie szerzej ujęte w kontekście analizy łańcucha podróży) i w konsekwencji drastycznie redukuje ostateczną dostępność do rynków celu~\cite{c-space, l-space2}.

\subsection{Inne reprezentacje}
Zaawansowana analityka sieci transportowych, wykraczając poza dualizm L-space i P-space, posiłkuje się w literaturze przedmiotu również modelami C-space oraz B-space. Topologia C-space (przestrzeń połączeń) odwraca klasyczną logikę grafową – węzłami stają się w niej linie komunikacyjne, połączone krawędzią w sytuacji, gdy współdzielą co najmniej jeden punkt przesiadkowy. Model B-space opiera się natomiast na grafach dwudzielnych (ang. \textit{bipartite graph}), tworząc dwa rozłączne zbiory węzłów (przystanki oraz linie), między którymi wyznacza się relacje~\cite{l-space2, c-space}.

Zastosowanie tych przestrzeni wykracza jednak poza ramy metodyczne niniejszego opracowania. Topologia C-space została pominięta, ponieważ badanie strukturalnych powiązań między samymi usługami nie jest głównym celem pracy. Podobnie odrzucono zastosowanie grafów dwudzielnych (B-space), rezygnując z mikroskopowego projektowania precyzyjnych macierzy przesiadek. Ze względu na dużą rozległość przestrzenną analizowanego obszaru, w to miejsce wdrożono silnie ugruntowane w analityce transportowej podejście makroskopowe. Opiera się ono na estymacji czasu na przesiadkę bazującej na uśrednionej częstotliwości kursowania docelowej linii, bez konieczności symulowania ruchu pojedynczych pojazdów~\cite{c-space, metropolitan, capetown}.


\section{Modele grawitacyjne w ocenie dostępności}\label{section:modele_grawitacyjne}

Modele grawitacyjne, nazywane również modelami potencjału, stanowią kluczowe narzędzie w ewaluacji systemów transportowych, pozwalając na przejście od prostych miar topologicznych do zaawansowanej analizy interakcji przestrzennych. Koncepcja ta, wywodząca się z pionierskich prac W. Hansena z 1959 roku, opiera się na założeniu, że dostępność danego obszaru jest wypadkową dwóch przeciwstawnych sił: atrakcyjności potencjalnych celów podróży oraz impedancji, czyli oporu przestrzeni utrudniającego dotarcie do nich~\cite{Koźlak_Aleksandra_Znaczenie_2010, gravity2, capetown}.

Aby rzetelnie ocenić użyteczność sieci, model musi uwzględniać wagę punktów docelowych. Ich atrakcyjność wyraża się najczęściej wymiernymi wskaźnikami społeczno-gospodarczymi, takimi jak całkowita liczba miejsc pracy, pojemność placówek zdrowotnych i edukacyjnych czy wielkość powierzchni handlowej. Z kolei wysiłek niezbędny do pokonania dystansu między strefą początkową a docelową modeluje funkcja impedancji. W badaniach nad zachowaniami przestrzennymi pasażerów funkcja ta przyjmuje zazwyczaj postać ujemnej funkcji wykładniczej lub potęgowej, co odzwierciedla spadek skłonności do podjęcia podróży wraz ze wzrostem jej uogólnionego kosztu. Kluczową rolę odgrywa tu wyznaczany empirycznie współczynnik zaniku (ang. \textit{distance decay parameter}), opisujący stopień wrażliwości danej grupy użytkowników na narastający opór przestrzeni~\cite{Koźlak_Aleksandra_Znaczenie_2010, gravity, gravity2, capetown}.

Pomimo dużej wartości analitycznej, klasyczne, nieograniczone modele grawitacyjne opierają się na uproszczonym założeniu o nieskończonej pojemności rynków docelowych. W gęstej tkance miejskiej prowadzi to do zignorowania zjawiska konkurencji o zasoby – pomija się fakt, że o ograniczoną pulę wakatów w centrum ubiegają się mieszkańcy z wielu, często bardzo odległych stref pochodzenia. Odpowiedzią na ten problem, stanowiącą obecnie standard w metodyce planowania przestrzennego, jest implementacja podwójnie ograniczonych modeli grawitacyjnych (ang. \textit{doubly constrained spatial interaction models}). Algorytmy te precyzyjnie bilansują przepływy: z jednej strony gwarantują zgodność z popytem generowanym w strefach zamieszkania, a z drugiej ściśle limitują napływ do stref docelowych w oparciu o ich faktyczną pojemność~\cite{gravity2}.


\section{Modelowanie całego łańcucha podróży (WTC)}\label{section:wtc}

Poprawne wyznaczenie funkcji impedancji, stanowiącej podstawę omówionych wcześniej modeli grawitacyjnych, wymaga precyzyjnego zdefiniowania szeroko rozumianego kosztu przemieszczania się pomiędzy strefami. W zaawansowanej analityce transportowej odchodzi się od uproszczonego podejścia, w którym opór przestrzeni utożsamiany jest wyłącznie z czasem spędzonym przez pasażera wewnątrz pojazdu. Z perspektywy rzetelnej ewaluacji systemów transportowych każda podróż jest wieloetapowym procesem, który należy modelować jako cały łańcuch podróży (ang. \textit{Whole Travel Chain} – WTC), realizując postulat pełnej mobilności „od drzwi do drzwi” (ang. \textit{door-to-door mobility})~\cite{disabilities1}.

\subsection{Pierwsza i ostatnia mila}\label{subsection:mila}
Kluczowym etapem budującym ten łańcuch jest uwzględnienie problemu pierwszej i ostatniej mili (ang. \textit{access/egress time}). Wynika to z faktu, że niemal każda podróż w publicznym transporcie zbiorowym rozpoczyna się i kończy ruchem pieszym lub z wykorzystaniem innych środków dowozowych. Czas dojścia z faktycznego punktu początkowego do przystanku startowego oraz czas przejścia z przystanku końcowego do ostatecznego celu podróży to parametry, które fizycznie limitują realną strefę oddziaływania węzła przesiadkowego. Ich pominięcie w modelach przestrzennych prowadzi do znacznego przeszacowania dostępności sieci. W analizach z wykorzystaniem narzędzi GIS czas ten wyznacza się najczęściej na podstawie odległości w zdefiniowanej sieci ciągów pieszych, podzielonej przez uśrednioną prędkość chodu~\cite{metropolitan, capetown}.

\subsection{Czas oczekiwania}\label{subsection:waiting}
Kolejnym istotnym elementem łańcucha podróży jest skumulowany czas oczekiwania (ang. \textit{waiting time}) na środek transportu, występujący zarówno przy inicjacji podróży, jak i w trakcie każdej przesiadki. W modelowaniu topologicznym opartym na rozkładowych bazach danych w formacie GTFS, przeciętny czas oczekiwania pasażera estymuje się powszechnie jako połowę interwału kursowania danej linii w analizowanym oknie czasowym. Warto podkreślić, że morfologia współczesnych obszarów metropolitalnych rzadko umożliwia obsłużenie wszystkich relacji popytowych wyłącznie za pomocą połączeń bezpośrednich, co naturalnie wymusza występowanie transferów wewnątrz systemu~\cite{LUO2019102505, metropolitan}.

\subsection{Penalizacja przesiadek}\label{subsection:przesiadki}
Proces przesiadania się stanowi dla użytkownika dodatkowe obciążenie. Poza fizyczną stratą czasu na przemieszczenie się między przystankami i oczekiwanie na kolejny pojazd, wprowadza on dyskomfort psychologiczny i wymaga nakładu dodatkowego wysiłku. W celu wiarygodnego odzwierciedlenia tego oporu w algorytmach wyznaczania najkrótszych ścieżek, stosuje się penalizację transferów. Sprowadza się to zazwyczaj do zaimplementowania tzw. kary przesiadkowej (ang. \textit{transfer penalty cost}), przybierającej postać ekwiwalentu czasu doliczanego do całkowitego kosztu podróży za każdą zmianę linii. W badaniach empirycznych nad sieciami transportowymi często przyjmuje się zryczałtowaną wartość karną (np. 5 minut dodatkowego obciążenia) za sam fakt wymuszenia na pasażerze transferu~\cite{LUO2019102505}.

\subsection{Uogólniony Koszt Podróży (GTC)}\label{subsection:gtc}
Agregacja omówionych wyżej składowych pozwala na wyznaczenie kompleksowej miary syntetycznej, jaką jest Uogólniony Koszt Podróży (ang. \textit{Generalized Travel Cost} – GTC). W analityce układów transportowych dopiero tak wyliczony koszt stanowi poprawną metodologicznie bazę do estymacji parametru impedancji w modelach grawitacyjnych, ponieważ obiektywnie odwzorowuje faktyczną barierę wysiłku pokonywania przestrzeni przez człowieka. Bez wdrożenia modelu całego łańcucha podróży i funkcji GTC do analiz interakcji, pomiary dostępności rynków pracy czy usług mogą prowadzić do błędnych wniosków, oderwanych od logiki funkcjonowania systemów operacyjnych i rzeczywistych zachowań przestrzennych podróżnych~\cite{LUO2019102505}.


\section{Sprawiedliwość transportowa i mapowanie społeczne}\label{section:sprawiedliwość}

Kompleksowa ocena publicznych systemów transportowych, wykraczająca poza ramy analityki topologicznej i wyznaczanie uogólnionego kosztu podróży, wymaga osadzenia uzyskanych wyników w kontekście struktury społecznej. W nowoczesnej inżynierii transportu i planowaniu przestrzennym istotne staje się rozróżnienie pomiędzy pojęciami równości (ang. \textit{equality}) a sprawiedliwości transportowej (ang. \textit{transport equity}). Równość definiowana jest najczęściej jako czysto geometryczne, równomierne rozmieszczenie infrastruktury lub zapewnienie identycznego poziomu usług dla ogółu populacji, co w praktyce często pomija zróżnicowanie fizycznych i ekonomicznych potrzeb mieszkańców. Sprawiedliwość transportowa poszukuje z kolei obiektywnej proporcjonalności w dystrybucji dostępu do mobilności, kładąc szczególny nacisk na systematyczne wsparcie społeczności niedostatecznie obsługiwanych przez sieć. Stanowi to wyraźne odejście od paradygmatu homogenicznego pasażera na rzecz uwzględnienia faktu, że obciążenia oraz koszty wynikające z przemieszczania się są wewnątrz tkanki miejskiej dystrybuowane asymetrycznie~\cite{KHALIFE2023160, elavating}.

Weryfikacja tego zjawiska na poziomie modelowania matematycznego wymusza zastosowanie segmentacji demograficznej. We wprowadzonych wcześniej podwójnie ograniczonych modelach grawitacyjnych kluczową rolę odgrywa współczynnik zaniku oporu przestrzennego. Empiryczne analizy przepływów wskazują, że wartość tego parametru wykazuje silną korelację ze statusem materialnym badanej grupy~\cite{capetown, gravity2}. Z tego względu konieczna staje się jego odrębna kalibracja w zależności od profilu dochodowego. Osoby o niższym statusie materialnym, dysponujące ograniczonym budżetem rozporządzalnym, wykazują wyższą wrażliwość na koszty finansowe podróży – co bezpośrednio wpisuje się w problematykę przystępności cenowej – zmuszając je nierzadko do akceptacji dłuższego czasu dojazdu w celu optymalizacji wydatków. Z kolei grupy zamożniejsze charakteryzują się odmienną krzywą popytu – ich wrażliwość na straty czasowe jest wyższa, co modyfikuje skłonność do pokonywania oporu przestrzeni~\cite{capetown}.

% \chapter{Przegląd Literatury i Podstawy Metodologiczne}\label{chapter:literatura}

% \section{Teoretyczne ramy dostępności: cztery komponenty dostępności}

% % Dostępność jest pojęciem wielowymiarowym, definiowanym w literaturze jako łatwość, z jaką mieszkańcy danego obszaru mogą dotrzeć do miejsc aktywności, takich jak miejsca pracy, usługi zdrowotne czy edukacyjne. W ujęciu teoretycznym wyróżnia się cztery kluczowe komponenty dostępności, które determinują jej poziom. Są to komponenty: transportowy, przestrzenny (użytkowania terenu), czasowy oraz indywidualny.

% % Komponent transportowy opisuje system transportu, uwzględniając czas podróży, koszty oraz wysiłek związany z przemieszczaniem się. Komponent przestrzenny (land-use) odnosi się do rozmieszczenia celów podróży (np. miejsc pracy) i ich atrakcyjności. Komponent czasowy uwzględnia ograniczenia czasowe, w jakich dostępne są usługi (np. godziny otwarcia) oraz czas, jakim dysponują użytkownicy. Wreszcie komponent indywidualny koncentruje się na potrzebach i możliwościach konkretnych użytkowników, takich jak dochód, wiek czy stopień sprawności fizycznej, co determinuje ich rzeczywistą zdolność do skorzystania z systemu transportowego.


% W paradygmacie zaawansowanej analizy przestrzennej i~inżynierii transportu, dostępność (ang.~accessibility) nie jest wyłącznie pochodną odległości fizycznej, lecz stanowi fundamentalny \textit{produkt} zintegrowanego systemu transportowego, decydujący o przewadze lokalizacyjnej danego obszaru~\cite{turoń_sierpiński_2019}. W~ujęciu formalnym, dostępność definiuje się jako stopień, w~jakim połączone systemy zagospodarowania przestrzennego i~transportu masowego umożliwiają określonym grupom jednostek dotarcie do pożądanych celów aktywności (takich jak miejsca pracy, edukacji czy szpitale) przy użyciu dostępnych środków przemieszczania się w~określonym czasie~\cite{capetown}. Aby dogłębnie diagnozować złożone interakcje sieciowe, ewaluacja dostępności opiera się na rygorystycznych ramach metodologicznych, które dekomponują to zjawisko na cztery współzależne komponenty: transportowy, przestrzenny, czasowy oraz indywidualny~\cite{gravity2, literature_review}.


% % \subsection{Komponent transportowy (podażowy i infrastrukturalny)}
% % Jest to rdzeń analizy sieciowej, określający wydajność, topologię i~koszty pokonywania przestrzeni przy użyciu systemu transportu~\cite{mrt}. Z~perspektywy teorii grafów i~nauki o sieciach, strukturę ta modelowana jest przy wykorzystaniu specyficznych topologii przestrzennych: L-space -- przestrzeń infrastruktury fizycznej, gdzie węzły-przystanki łączone są krawędziami drogowymi/torowymi, P-space -- przestrzeń usług, ilustrująca bezpośrednie połączenia bezpośrednie (bez przesiadek pomiędzy liniami), a także C-space, gdzie węzłami są same linie komunikacyjne, co doskonale kwantyfikuje nakładanie się na siebie usług transportowych i konieczność transferów~\cite{l-space2, c-space, LUO2019102505}. Z punktu widzenia pasażera, komponent transportowy reprezentowany jest przez funkcję impedancji (oporu przestrzeni), która agreguje czas dojścia, czas oczekiwania, czas jazdy oraz uciążliwość przesiadek w jedną wartość, tzw. Uogólniony Koszt Podróży (GTC – Generalized Travel Cost). W zaawansowanym modelowaniu, każda przesiadka obarczona jest dodatkową karą czasową (tzw. penalty cost), drastycznie obniżającą ogólną dostępność topologiczną danej relacji wewnątrz sieci~\cite{capetown, LUO2019102505}.

% \subsection{Komponent transportowy (podażowy i infrastrukturalny)}

% Komponent ten stanowi rdzeń podażowy całego systemu przemieszczania się, definiując fizyczne oraz operacyjne możliwości pokonywania odległości pomiędzy strefami źródłowymi a docelowymi. W najszerszym ujęciu odzwierciedla on wielkość oporu przestrzeni, z jakim musi zmierzyć się użytkownik sieci transportowej. Na wymiar ten składają się dwa nierozerwalnie połączone aspekty: twarda infrastruktura (układ dróg, torowisk, lokalizacja przystanków i węzłów przesiadkowych) oraz nałożona na nią miękka warstwa operacyjna (przebieg linii komunikacyjnych, parametry taboru, prędkości eksploatacyjne czy pojemność systemu)~\cite{mrt, l-space2, c-space}.

% Z perspektywy analitycznej, komponent ten pełni rolę głównego pośrednika determinującego wydajność obsługi obszaru miejskiego. Jego rygorystyczna ocena prowadzona jest dwutorowo. Z jednej strony, strukturę i układ relacji wewnątrz systemu bada się przy użyciu aparatury matematycznej wywodzącej się z teorii grafów, co pozwala na identyfikację wąskich gardeł i ewaluację spójności powiązań. Z drugiej strony, z punktu widzenia pasażera, organizacja transportu bezpośrednio przekłada się na całkowity wysiłek i czas niezbędny do pokonania trasy. Wymusza to syntetyczną kwantyfikację tzw. impedancji, czyli uogólnionego kosztu podróży, który uwzględnia wszystkie obciążenia związane z przemieszczaniem się. W konsekwencji komponent transportowy nie tylko warunkuje fizyczną możliwość dotarcia do celu, ale poprzez uciążliwość samej podróży bezpośrednio kształtuje postrzeganą dostępność i atrakcyjność danego obszaru~\cite{capetown, LUO2019102505}.

% % \subsection{Komponent przestrzenny (land-use)}
% % Komponent ten definiuje masę i rozmieszczenie rynków celu w przestrzeni miejskiej (np. lokalizację miejsc pracy, handlu, ochrony zdrowia)~\cite{elavating, capetown, Koźlak_Aleksandra_Znaczenie_2010}. W nowoczesnej analizie przestrzennej nie ograniczamy się do prostego zliczania skumulowanych szans (ang. cumulative opportunities). Kluczowe jest wdrożenie podwójnie ograniczonych modeli grawitacyjnych (ang. doubly constrained spatial interaction models), które uwzględniają zarówno podaż, przepustowość w punkcie docelowym, jak i popyt generowany w strefach pochodzenia. Innymi słowy, ocena dostępności musi uwzględniać efekt konkurencji przestrzennej – atrakcyjność danego klastra miejsc pracy maleje proporcjonalnie do liczby innych użytkowników sieci rywalizujących o te same zasoby i wakaty~\cite{gravity2, mrt}.


% % \subsection{Komponent czasowy}
% % Dostępność w transporcie zbiorowym wykazuje potężną fluktuację w układzie czasoprzestrzennym. Komponent czasowy uwzględnia ograniczenia związane z godzinami otwarcia celów podróży, harmonogramami pracy oraz, co najważniejsze, rozkładami jazdy determinującymi częstotliwość i niezawodność usług przewozowych~\cite{gravity2, capetown, review, LUO2019102505}. Do integracji tego wymiaru w sieciach złożonych standardowo wykorzystujemy relacyjne bazy danych GTFS (General Transit Feed Specification), które pozwalają na wyznaczanie dynamicznych czasów podróży oraz uwzględniają czas oczekiwania jako połowę interwału kursowania (ang. headway) między pojazdami. Analiza w tym komponencie dowodzi, że obszary o rzekomo doskonałej dostępności w godzinach szczytu mogą cierpieć na drastyczne wykluczenie transportowe wieczorami, co zaburza spójność sieci~\cite{capetown, LUO2019102505, elavating}.


% % \subsection{Komponent indywidualny (socjo-ekonomiczny i fizyczny)}
% % Dostępność sieciowa ma charakter silnie zindywidualizowany – ta sama topologia transportowa oferuje diametralnie różny poziom dostępu do celów w zależności od możliwości danego użytkownika~\cite{gravity2, capetown, mrt}. Ten komponent nakazuje badanie wrażliwości różnych grup demograficznych. Oprócz barier fizycznych i architektonicznych, które znacząco ograniczają mobilność osób z niepełnosprawnościami, w ujęciu systemowym absolutnie kluczowa staje się przystępność cenowa usług (ang. affordability)~\cite{capetown, elavating}. Z perspektywy analitycznej oznacza to, że nawet najlepiej zaprojektowana infrastruktura staje się bezużyteczna, jeśli cena przejazdu przewyższa możliwości finansowe mieszkańców. Z tego względu nowoczesne narzędzia analizy przestrzennej muszą integrować progi przystępności cenowej z funkcjami oporu przestrzennego (impedancji) oddzielnie dla każdej grupy dochodowej~\cite{capetown}.


% % Z perspektywy analizy sieciowej, system ten diagnozujemy poprzez dualizm topologiczny: L-space (topologia infrastruktury fizycznej, np. torowisk i dróg) oraz P-space (topologia samych usług, w której łączymy węzły bezprzesiadkowymi relacjami, co doskonale obrazuje logikę podróży pasażera)~\cite{LUO2019102505}. Komponent transportowy kwantyfikuje wysiłek niezbędny do pokonania odległości, co matematycznie wyraża funkcja impedancji (oporu przestrzeni)~\cite{capetown, Koźlak_Aleksandra_Znaczenie_2010}. Dla transportu zbiorowego ten „opór” to Uogólniony Koszt Podróży (GTC – Generalized Travel Cost), który rygorystycznie agreguje czas dojścia pieszego do przystanku, czas oczekiwania (zależny od częstotliwości), czas jazdy oraz – co krytyczne w analizie grafów przesiadkowych – uciążliwość i czas samych transferów, w tym stosowane w modelach „kary czasowe” (penalty cost) za każdą zmianę środka transportu~\cite{LUO2019102505, literature_review}.

% \subsection{Komponent przestrzenny (land-use)}
% Z drugiej strony równania dostępności znajduje się wymiar przestrzenny. Odnosi się on bezpośrednio do rozmieszczenia, wielkości oraz zróżnicowania rynków celu w tkance miejskiej (takich jak miejsca pracy, handlu, ochrony zdrowia czy edukacji)~\cite{elavating, capetown, Koźlak_Aleksandra_Znaczenie_2010}. W modelach transportowych obszar ten reprezentuje miarę atrakcyjności poszczególnych stref, stanowiąc docelowy punkt podróży użytkowników sieci. Współczesne ujęcie tego zagadnienia odchodzi od traktowania celów wyłącznie jako izolowanych punktów na mapie, zwracając uwagę na ograniczoną pojemność tych miejsc oraz zjawisko wzajemnej konkurencji o zasoby (np. wakaty na rynku pracy). Precyzyjne bilansowanie popytu z faktycznymi możliwościami rynków docelowych stanowi wyzwanie, które w zaawansowanej analityce rozwiązuje się poprzez implementację odpowiednich modeli grawitacyjnych, co zostało szerzej omówione w dalszej części pracy~\cite{gravity2, mrt}.

% \subsection{Komponent czasowy}
% Oprócz fizycznego pokonywania odległości, krytycznym aspektem ewaluacji systemów transportowych jest ich zmienność. Dostępność, ze szczególnym uwzględnieniem komunikacji zbiorowej, rzadko bywa wartością stałą – wykazuje za to potężną fluktuację w zależności od pory dnia czy tygodnia. Wymiar czasowy spina ze sobą ograniczenia nakładane z jednej strony przez harmonogramy funkcjonowania samych celów podróży (np. godziny otwarcia placówek czy ramy czasu pracy), a z drugiej – przez rozkłady jazdy, które bezpośrednio determinują częstotliwość i niezawodność usług przewozowych~\cite{gravity2, capetown, review, LUO2019102505}. Zrozumienie tej dynamiki uświadamia, że obszary o rzekomo doskonałej łączności w godzinach szczytu mogą cierpieć na drastyczne wykluczenie wieczorami. Rygorystyczne ujęcie tych wahań wymaga włączenia do analizy czasu oczekiwania na przystanku, co zostanie rozwinięte w rozdziale poświęconym modelowaniu całego łańcucha podróży~\cite{capetown, LUO2019102505, elavating}.

% \subsection{Komponent indywidualny (socjo-ekonomiczny i fizyczny)}
% Całościowy obraz analizowanego zjawiska domyka perspektywa samego pasażera. Należy bowiem podkreślić, że obiektywnie ta sama sieć może oferować diametralnie różny poziom możliwości przemieszczania się w zależności od cech i ograniczeń konkretnego użytkownika~\cite{gravity2, capetown, mrt}. Aspekt ten nakazuje ewaluację systemu przez pryzmat barier napotykanych przez poszczególne grupy demograficzne. Obejmuje to zarówno bariery fizyczne i architektoniczne, utrudniające podróżowanie osobom z niepełnosprawnościami, jak i uwarunkowania socjoekonomiczne, wśród których kluczowa staje się przystępność cenowa usług~\cite{capetown, elavating}. Z perspektywy systemowej, nawet najlepiej zaprojektowana infrastruktura staje się bezużyteczna, jeśli koszt przejazdu przewyższa budżet mieszkańców. Zjawisko to i wynikająca z niego konieczność segmentacji podróżnych stanowią fundament zagadnień sprawiedliwości transportowej, poruszanych w kolejnych sekcjach~\cite{capetown}.


% \section{Topologiczna analiza sieci}

% Rygorystyczna ocena wydajności transportu zbiorowego – stanowiąca sedno komponentu transportowego – wymaga mapowania rzeczywistych układów przy użyciu matematycznej teorii grafów. Nowoczesne podejście do modelowania złożonych systemów transportowych rozwiązuje ten problem, dekomponując ich strukturę na odrębne, lecz ściśle powiązane warstwy topologiczne: L-Space oraz P-Space. Zrozumienie tego dualizmu (oddzielenia fizycznej infrastruktury od logiki usług) stanowi kluczowy warunek w inżynierii transportu do identyfikowania wąskich gardeł i błędów operacyjnych w zarządzaniu ruchem~\cite{l-space2}.

% % Jak zasygnalizowano przy okazji definicji komponentu transportowego, rygorystyczna ocena wydajności systemu publicznego wymaga zaawansowanego modelowania jego struktury za pomocą teorii grafów. W nurcie badawczym analizy sieci złożonych wykształcił się ustandaryzowany aparat pojęciowy, w ramach którego strukturę transportu zbiorowego dekomponuje się na powiązane ze sobą warstwy topologiczne, a w szczególności modele L-Space oraz P-Space. Z perspektywy zaawansowanej inżynierii transportu, pełne zrozumienie tego dualizmu jest absolutnym fundamentem pozwalającym diagnozować wady infrastrukturalne oraz operacyjne układów komunikacyjnych~\cite{l-space2}.

% % \subsection{Przestrzeń infrastruktury (L-Space)}
% % Topologia L-space to reprezentacja sieciowa, w której węzłami są fizyczne przystanki lub stacje, a krawędź między nimi tworzona jest wtedy i tylko wtedy, gdy dwa przystanki następują bezpośrednio po sobie na trasie i łączy je fizycznie co najmniej jedna linia komunikacyjna. Model ten stanowi surowe odwzorowanie fizycznej infrastruktury miejskiej – układu szlaków, torowisk czy ulic nałożonych na topografię terenu. Z perspektywy analitycznej, badanie L-space jest niezastąpione w diagnozowaniu problemów związanych stricte z podażą: pozwala na precyzyjną identyfikację wąskich gardeł (bottlenecks), określanie bezwzględnych odległości przestrzennych oddzielających węzły oraz badanie fizycznej odporności (ang. vulnerability) całego układu na awarie czy zablokowanie konkretnego odcinka~\cite{c-space, l-space2}.

% % \subsection{Przestrzeń usług (P-space)}
% % Całkowicie odmienną, zorientowaną na użytkownika perspektywę oferuje topologia P-space (przestrzeń usług). W tej przestrzeni węzłami wprawdzie nadal pozostają przystanki, jednak krawędź łączy wybrane dwa węzły wtedy, gdy mogą one zostać osiągnięte bezpośrednio, to jest poprzez podróż pojedynczą linią, bez jakiejkolwiek konieczności przesiadania się. Skutkuje to doniosłą z punktu widzenia matematyki i logistyki własnością: w P-space wszystkie przystanki obsługiwane przez konkretną linię tworzą pełny podgraf, klikę. Jest to reprezentacja wprost uwypuklająca to, jak pasażer operacyjnie postrzega system. W tym ujęciu odległość topologiczna, czyli najkrótsza ścieżka między wierzchołkami w grafie, bezpośrednio pozwala kwantyfikować minimalną liczbę linii niezbędnych do dotarcia do celu~\cite{c-space, l-space2}. 

% \subsection{Przestrzeń infrastruktury (L-Space)}
% Topologia L-space stanowi podstawową reprezentację sieciową, w której węzły odpowiadają fizycznym przystankom lub stacjom, natomiast łączące je krawędzie definiuje się wyłącznie między przystankami następującymi bezpośrednio po sobie na trasie danej linii komunikacyjnej. Model ten jest w istocie wiernym odwzorowaniem fizycznej infrastruktury miejskiej – układu ulic, torowisk i szlaków osadzonych w topografii terenu. Z perspektywy inżynierii transportu badanie przestrzeni L-space okazuje się niezastąpione w diagnozowaniu problemów o charakterze podażowym. Pozwala ono na precyzyjną identyfikację wąskich gardeł (ang. \textit{bottlenecks}), estymację rzeczywistych odległości przestrzennych oddzielających węzły oraz ocenę fizycznej odporności (ang. \textit{vulnerability}) całego układu na awarie i zakłócenia w ruchu~\cite{c-space, l-space2}.

% \subsection{Przestrzeń usług (P-space)}
% Z kolei topologia P-space oferuje całkowicie odmienną, zorientowaną na pasażera perspektywę analityczną. Chociaż zbiór węzłów nadal opiera się na fizycznych przystankach, to krawędzie łączą w niej wybrane strefy wtedy, gdy można między nimi podróżować bezpośrednio – przy użyciu pojedynczej linii komunikacyjnej, bez konieczności przesiadania się. Taka konstrukcja prowadzi do kluczowej własności matematycznej: w modelu P-space wszystkie przystanki obsługiwane przez konkretną linię tworzą pełny podgraf (klikę). Reprezentacja ta doskonale odzwierciedla operacyjne postrzeganie układu przez użytkownika. W tym ujęciu odległość topologiczna, czyli najkrótsza ścieżka wyznaczona w grafie, nie określa dystansu w kilometrach, lecz bezpośrednio kwantyfikuje minimalną liczbę linii niezbędnych do osiągnięcia celu podróży~\cite{c-space, l-space2}.

% % \subsection{Istotność separacji infrastruktury i usług}

% % Znaczenie analitycznego dualizmu L-space/P-space polega na ścisłym oddzieleniu infrastruktury (tego, co fizycznie istnieje w przestrzeni) od nałożonego na nią schematu operacyjnego (tego, jak ta przestrzeń jest wykorzystywana poprzez siatkę połączeń). Dopiero zestawienie ze sobą wskaźników sieciowych z obu tych topologii pozwala ujawnić utajone dysfunkcje układu. 

% % Przykładowo, jeżeli dany system transportowy charakteryzuje się wysoką spójnością w L-space, ale wykazuje duże średnie długości ścieżek w P-space, implikuje to błędny model zarządzania operacyjnego. Oznacza to bowiem, że pomimo doskonałej fizycznej siatki ulic czy torów, organizator narzucił na nią wysoce pofragmentowane usługi (np. krótkie trasy linii), zmuszając pasażerów do ciągłych przesiadek. Jak nadmieniono wcześniej, każda przesiadka winduje w górę Uogólniony Koszt Podróży (GTC) poprzez tzw. karę czasową (penalty cost), drastycznie tnąc w dół ostateczną dostępność do rynków celu~\cite{c-space, l-space2}. 

% \subsection{Istotność separacji infrastruktury i usług}

% Znaczenie analitycznego dualizmu L-space/P-space polega na ścisłym oddzieleniu infrastruktury (tego, co fizycznie istnieje w przestrzeni) od nałożonego na nią schematu operacyjnego (tego, jak owa przestrzeń jest wykorzystywana poprzez siatkę połączeń). Dopiero zestawienie ze sobą wskaźników sieciowych z obu tych topologii pozwala ujawnić utajone dysfunkcje układu.

% Przykładowo, jeżeli dany system transportowy charakteryzuje się wysoką spójnością w topologii L-space, ale jednocześnie wykazuje duże średnie długości ścieżek w przestrzeni P-space, implikuje to błędny model zarządzania operacyjnego. Oznacza to bowiem, że pomimo doskonałej fizycznej siatki ulic czy torów, organizator narzucił na nią wysoce pofragmentowane usługi (np. bardzo krótkie trasy linii), zmuszając pasażerów do ciągłych przesiadek. Taka sytuacja drastycznie zwiększa uciążliwość przemieszczania się – każdy wymuszony transfer stanowi istotną barierę dla pasażera, która winduje całkowity wysiłek wkładany w podróż (co zostanie szerzej ujęte w kontekście analizy łańcucha podróży) i w konsekwencji drastycznie redukuje ostateczną dostępność do rynków celu~\cite{c-space, l-space2}.

% % \subsection{Inne reprezentacje}

% % Dopełnieniem topologicznej analityki sieci transportowych, wykraczającym poza opisany wcześniej dualizm, jest w literaturze przedmiotu wdrożenie modeli C-space oraz B-space. Topologia C-space (przestrzeń tras/usług) odwraca klasyczną logikę grafową: węzłami stają się w niej same linie komunikacyjne, a krawędzie definiuje się na podstawie współdzielenia przez nie co najmniej jednego punktu transferowego (przystanku). Z kolei model B-space opiera się na zastosowaniu grafów dwudzielnych (ang. bipartite graph), w których ustanawia się dwa całkowicie rozłączne zbiory węzłów – fizyczne przystanki oraz linie transportowe – łącząc je krawędziami tylko pomiędzy tymi dwoma zbiorami~\cite{l-space2, c-space}. 

% % Należy jednak zaznaczyć, że pod kątem niniejszej pracy powyższe przestrzenie nie są istotne. W przypadku topologii C-space wynika to z faktu, że w celem poniższych analiz nie jest badanie interakcji pomiędzy poszczególnymi połączeniami. Z kolei w odniesieniu do topologii B-space jej badanie zostało pominięte, ponieważ nie stanowi przedmiotu rozważań mikroskopowe projektowanie macierzy przesiadek. Zamiast tego, mając na uwadze znaczną rozległość przestrzenną prowadzonych analiz sieciowych, zdecydowano się na implementację powszechnie uznanego i silnie ugruntowanego w literaturze przedmiotu uproszczenia. Polega ono na estymacji czasu oczekiwania na transfer na podstawie średniego interwału kursowania danej linii, gdzie zgodnie ze standardami analityki transportowej najczęściej przyjmuje się wartość równą połowie tegoż interwału~\cite{c-space, metropolitan, capetown}.

% \subsection{Inne reprezentacje}

% Zaawansowana analityka sieci transportowych, wykraczając poza dualizm L-space i P-space, posiłkuje się w literaturze przedmiotu również modelami C-space oraz B-space. Topologia C-space (przestrzeń połączeń) odwraca klasyczną logikę grafową – węzłami stają się w niej linie komunikacyjne, połączone krawędzią w sytuacji, gdy współdzielą co najmniej jeden punkt przesiadkowy. Model B-space opiera się natomiast na grafach dwudzielnych (ang. \textit{bipartite graph}), tworząc dwa rozłączne zbiory węzłów (przystanki oraz linie), między którymi wyznacza się relacje~\cite{l-space2, c-space}.

% Zastosowanie tych przestrzeni wykracza jednak poza ramy metodyczne niniejszego opracowania. Topologia C-space została pominięta, ponieważ badanie strukturalnych powiązań między samymi usługami nie jest głównym celem pracy. Podobnie odrzucono zastosowanie grafów dwudzielnych (B-space), rezygnując z mikroskopowego projektowania precyzyjnych macierzy przesiadek. Ze względu na dużą rozległość przestrzenną analizowanego obszaru, w to miejsce wdrożono silnie ugruntowane w analityce transportowej podejście makroskopowe. Opiera się ono na estymacji czasu oczekiwania na przesiadkę za pomocą miary probabilistycznej – przyjmując, że uśredniony czas oczekiwania na przystanku odpowiada połowie interwału kursowania docelowej linii komunikacyjnej~\cite{c-space, metropolitan, capetown}.

% % \section{Modele grawitacyjne w ocenie dostępności}

% % Modele grawitacyjne, nazywane również modelami potencjału, stanowią absolutny fundament ewaluacji złożonych systemów transportowych, pozwalając na przejście od prostych miar topologicznych do rygorystycznej analityki interakcji przestrzennych. Koncepcja ta, wywodząca się z pionierskich prac W. Hansena z 1959 roku, opiera się na założeniu, że dostępność danego obszaru jest wypadkową dwóch przeciwstawnych sił: atrakcyjności potencjalnych celów podróży oraz impedancji, czyli oporu przestrzeni utrudniającego dotarcie do nich~\cite{Koźlak_Aleksandra_Znaczenie_2010, gravity2, capetown}.

% % Aby rzetelnie ocenić użyteczność sieci, model musi uwzględniać to, jak bardzo pożądane są punkty docelowe. Ich atrakcyjność wyraża się najczęściej konkretnymi miarami społeczno-gospodarczymi, takimi jak całkowita liczba miejsc pracy, pojemność placówek zdrowotnych, edukacyjnych czy wielkość powierzchni handlowej. Z kolei wysiłek niezbędny do pokonania dystansu między strefą początkową a docelową modeluje funkcja impedancji. W badaniach nad przestrzennymi zachowaniami pasażerów funkcja ta przyjmuje najczęściej postać ujemnej funkcji wykładniczej lub potęgowej, co odzwierciedla spadek skłonności do podjęcia podróży wraz ze wzrostem jej uogólnionego kosztu. Kluczową rolę odgrywa tu wyznaczany empirycznie współczynnik zaniku, precyzyjnie opisujący stopień wrażliwości danej grupy użytkowników na narastający opór przestrzeni~\cite{Koźlak_Aleksandra_Znaczenie_2010, gravity, gravity2, capetown}.

% % Mimo ogromnej wartości analitycznej, klasyczne, nieograniczone modele grawitacyjne opierają się na uproszczonym założeniu o nieskończonej pojemności rynków docelowych. W gęstej tkance miejskiej prowadzi to do ignorowania zjawiska konkurencji o zasoby – pomija się fakt, że o ograniczoną pulę wakatów w centrum ubiegają się mieszkańcy z wielu różnych, często odległych stref pochodzenia. Odpowiedzią na ten problem, stanowiącą obecnie standard w zaawansowanej analityce przestrzennej, jest wdrożenie podwójnie ograniczonych modeli grawitacyjnych. Algorytmy te precyzyjnie bilansują przepływy: z jednej strony gwarantując zgodność z popytem generowanym w strefach zamieszkania, a z drugiej bezwzględnie limitując napływ do stref docelowych w oparciu o ich faktyczną pojemność~\cite{gravity2}.

% \section{Modele grawitacyjne w ocenie dostępności}

% Modele grawitacyjne, nazywane również modelami potencjału, stanowią kluczowe narzędzie w ewaluacji systemów transportowych, pozwalając na przejście od prostych miar topologicznych do zaawansowanej analizy interakcji przestrzennych. Koncepcja ta, wywodząca się z pionierskich prac W. Hansena z 1959 roku, opiera się na założeniu, że dostępność danego obszaru jest wypadkową dwóch przeciwstawnych sił: atrakcyjności potencjalnych celów podróży oraz impedancji, czyli oporu przestrzeni utrudniającego dotarcie do nich~\cite{Koźlak_Aleksandra_Znaczenie_2010, gravity2, capetown}.

% Aby rzetelnie ocenić użyteczność sieci, model musi uwzględniać wagę punktów docelowych. Ich atrakcyjność wyraża się najczęściej wymiernymi wskaźnikami społeczno-gospodarczymi, takimi jak całkowita liczba miejsc pracy, pojemność placówek zdrowotnych i edukacyjnych czy wielkość powierzchni handlowej. Z kolei wysiłek niezbędny do pokonania dystansu między strefą początkową a docelową modeluje funkcja impedancji. W badaniach nad zachowaniami przestrzennymi pasażerów funkcja ta przyjmuje zazwyczaj postać ujemnej funkcji wykładniczej lub potęgowej, co odzwierciedla spadek skłonności do podjęcia podróży wraz ze wzrostem jej uogólnionego kosztu. Kluczową rolę odgrywa tu wyznaczany empirycznie współczynnik zaniku (ang. \textit{distance decay parameter}), opisujący stopień wrażliwości danej grupy użytkowników na narastający opór przestrzeni~\cite{Koźlak_Aleksandra_Znaczenie_2010, gravity, gravity2, capetown}.

% Pomimo dużej wartości analitycznej, klasyczne, nieograniczone modele grawitacyjne opierają się na uproszczonym założeniu o nieskończonej pojemności rynków docelowych. W gęstej tkance miejskiej prowadzi to do zignorowania zjawiska konkurencji o zasoby – pomija się fakt, że o ograniczoną pulę wakatów w centrum ubiegają się mieszkańcy z wielu, często bardzo odległych stref pochodzenia. Odpowiedzią na ten problem, stanowiącą obecnie standard w metodyce planowania przestrzennego, jest implementacja podwójnie ograniczonych modeli grawitacyjnych (ang. \textit{doubly constrained spatial interaction models}). Algorytmy te precyzyjnie bilansują przepływy: z jednej strony gwarantują zgodność z popytem generowanym w strefach zamieszkania, a z drugiej ściśle limitują napływ do stref docelowych w oparciu o ich faktyczną pojemność~\cite{gravity2}.

% % Jednym z najczęściej stosowanych narzędzi do oceny dostępności są modele potencjału oparte na prawie grawitacji, wprowadzone do planowania przez Hansena. Podstawowe założenie tych modeli głosi, że interakcja (dostępność $A_i$) między strefą pochodzenia $i$ a strefą celu $j$ jest wprost proporcjonalna do atrakcyjności celu ($T_{attr}$ lub $O_j$) i odwrotnie proporcjonalna do impedancji podróży.

% % Wzór ogólny przyjmuje postać:
% % $$A_i = \sum_j O_j f(c_{ij})$$
% % gdzie $O_j$ reprezentuje masę lub atrakcyjność celu (np. liczbę miejsc pracy, populację, powierzchnię handlową), a $f(c_{ij})$ jest funkcją impedancji (oporu przestrzeni). Funkcja impedancji najczęściej przybiera formę wykładniczą ($e^{-\beta c_{ij}}$) lub potęgową, odzwierciedlając spadek chęci do podróży wraz ze wzrostem kosztu lub czasu $c_{ij}$.

% % \section{Modelowanie całego łańcucha podróży (WTC)}

% % Poprawne wyznaczenie funkcji impedancji, będącej absolutnym fundamentem omówionych wcześniej modeli grawitacyjnych, wymaga precyzyjnego i rygorystycznego zdefiniowania szeroko rozumianego kosztu przemieszczania się pomiędzy strefami. W zaawansowanej analityce inżynierii transportu kategorycznie odrzuca się uproszczone podejście, w którym opór przestrzeni utożsamiany jest wyłącznie z czasem spędzonym przez pasażera wewnątrz pojedynczego pojazdu. Z perspektywy rzetelnej ewaluacji systemów transportowych każda podróż jest wieloetapowym, złożonym procesem, który należy modelować jako Cały Łańcuch Podróży (ang. Whole Travel Chain – WTC), realizując niezbędny w ocenie dostępności postulat pełnej mobilności „od drzwi do drzwi” (ang. door-to-door mobility)~\cite{disabilities1}.

% % \subsection{Pierwsza i ostatnia mila}

% % Krytycznym komponentem budującym ten łańcuch jest uwzględnienie problemu pierwszej i ostatniej mili (ang. access/egress), co wynika wprost z faktu, że absolutnie każda podróż w publicznym transporcie zbiorowym musi rozpocząć się i zakończyć ruchem pieszym lub z wykorzystaniem innych środków dowozowych. Czas dojścia z faktycznego punktu początkowego do przystanku startowego (access time) oraz czas dojścia z przystanku końcowego do ostatecznego celu podróży (egress time) to kluczowe parametry, które fizycznie limitują realną strefę oddziaływania węzła przesiadkowego, a ich zignorowanie skutkuje fałszywym przeszacowaniem modelowanej dostępności sieci. W badaniach przestrzennych z wykorzystaniem narzędzi GIS czas ten wyznacza się najczęściej na podstawie odległości w rygorystycznie zdefiniowanej sieci ciągów pieszych, podzielonej przez uśrednioną prędkość poruszania się~\cite{metropolitan, capetown}.

% % \subsection{Czas oczekiwania}

% % Kolejnym, istotnym elementem w łańcuchu podróży jest skumulowany czas oczekiwania (ang. waiting time) na środek transportu, występujący zarówno przy inicjacji podróży, jak i w trakcie każdej koniecznej przesiadki. W nowoczesnym modelowaniu topologicznym i grafowym, opartym na rozkładowych bazach danych w formacie GTFS (General Transit Feed Specification), przeciętny czas oczekiwania pasażera estymuje się powszechnie jako połowę interwału kursowania danej linii w analizowanym oknie czasowym. Należy przy tym podkreślić, że morfologia współczesnych obszarów metropolitalnych rzadko umożliwia obsłużenie wszystkich relacji popytowych za pomocą połączeń bezpośrednich, naturalnie wymuszając występowanie transferów wewnątrz systemu~\cite{LUO2019102505, metropolitan}.

% % \subsection{Penalizacja przesiadek}

% % Proces przesiadania się stanowi uciążliwe obciążenie dla użytkownika, ponieważ poza fizyczną stratą czasu na przemieszczenie się między przystankami lub peronami i oczekiwanie na kolejny pojazd, wprowadza on silny dyskomfort psychologiczny i wymaga dodatkowego wysiłku. W celu wiarygodnego odzwierciedlenia tego oporu w matematycznych algorytmach wyznaczania najkrótszych ścieżek, stosuje się bezwzględną penalizację transferów. Sprowadza się to zazwyczaj do zaimplementowania tak zwanej kary przesiadkowej (ang. transfer penalty cost), przybierającej postać ekwiwalentu czasu doliczanego do całkowitego czasu podróży za każdą dokonaną zmianę linii. W badaniach empirycznych nad sieciami transportowymi przyjmuje się nierzadko zryczałtowaną wartość karną, wynoszącą na przykład 5 minut dodatkowego obciążenia w funkcji kosztu za sam fakt wymuszenia na pasażerze transferu~\cite{LUO2019102505}.

% % \subsection{Uogólniony Koszt Podróży (GTC)}

% % Rygorystyczna agregacja omówionych wyżej składowych pozwala na wyznaczenie najbardziej pożądanej, syntetycznej miary, jaką jest Uogólniony Koszt Podróży (ang. Generalized Travel Cost – GTC). W świetle logistyki i analityki układów złożonych, dopiero tak kompleksowo wyliczony koszt stanowi właściwą i poprawną metodologicznie bazę do estymacji parametru impedancji w modelach grawitacyjnych, gdyż tylko on obiektywnie odwzorowuje faktyczną barierę wysiłku pokonywania przestrzeni przez człowieka. Bez rzetelnego wdrożenia modelu całego łańcucha podróży i funkcji GTC do analiz interakcji, pomiary dostępności rynków pracy czy usług pozostają jedynie teoretyczną iluzją, oderwaną od logiki funkcjonowania systemów operacyjnych i zachowań przestrzennych podróżnych~\cite{LUO2019102505}.

% \section{Modelowanie całego łańcucha podróży (WTC)}

% Poprawne wyznaczenie funkcji impedancji, stanowiącej podstawę omówionych wcześniej modeli grawitacyjnych, wymaga precyzyjnego zdefiniowania szeroko rozumianego kosztu przemieszczania się pomiędzy strefami. W zaawansowanej analityce transportowej odchodzi się od uproszczonego podejścia, w którym opór przestrzeni utożsamiany jest wyłącznie z czasem spędzonym przez pasażera wewnątrz pojazdu. Z perspektywy rzetelnej ewaluacji systemów transportowych każda podróż jest wieloetapowym procesem, który należy modelować jako cały łańcuch podróży (ang. \textit{Whole Travel Chain} – WTC), realizując postulat pełnej mobilności „od drzwi do drzwi” (ang. \textit{door-to-door mobility})~\cite{disabilities1}.

% \subsection{Pierwsza i ostatnia mila}

% Kluczowym etapem budującym ten łańcuch jest uwzględnienie problemu pierwszej i ostatniej mili (ang. \textit{access/egress time}). Wynika to z faktu, że niemal każda podróż w publicznym transporcie zbiorowym rozpoczyna się i kończy ruchem pieszym lub z wykorzystaniem innych środków dowozowych. Czas dojścia z faktycznego punktu początkowego do przystanku startowego oraz czas przejścia z przystanku końcowego do ostatecznego celu podróży to parametry, które fizycznie limitują realną strefę oddziaływania węzła przesiadkowego. Ich pominięcie w modelach przestrzennych prowadzi do znacznego przeszacowania dostępności sieci. W analizach z wykorzystaniem narzędzi GIS czas ten wyznacza się najczęściej na podstawie odległości w zdefiniowanej sieci ciągów pieszych, podzielonej przez uśrednioną prędkość chodu~\cite{metropolitan, capetown}.

% \subsection{Czas oczekiwania}

% Kolejnym istotnym elementem łańcucha podróży jest skumulowany czas oczekiwania (ang. \textit{waiting time}) na środek transportu, występujący zarówno przy inicjacji podróży, jak i w trakcie każdej przesiadki. W modelowaniu topologicznym opartym na rozkładowych bazach danych w formacie GTFS, przeciętny czas oczekiwania pasażera estymuje się powszechnie jako połowę interwału kursowania danej linii w analizowanym oknie czasowym. Warto podkreślić, że morfologia współczesnych obszarów metropolitalnych rzadko umożliwia obsłużenie wszystkich relacji popytowych wyłącznie za pomocą połączeń bezpośrednich, co naturalnie wymusza występowanie transferów wewnątrz systemu~\cite{LUO2019102505, metropolitan}.

% \subsection{Penalizacja przesiadek}

% Proces przesiadania się stanowi dla użytkownika dodatkowe obciążenie. Poza fizyczną stratą czasu na przemieszczenie się między przystankami i oczekiwanie na kolejny pojazd, wprowadza on dyskomfort psychologiczny i wymaga nakładu dodatkowego wysiłku. W celu wiarygodnego odzwierciedlenia tego oporu w algorytmach wyznaczania najkrótszych ścieżek, stosuje się penalizację transferów. Sprowadza się to zazwyczaj do zaimplementowania tzw. kary przesiadkowej (ang. \textit{transfer penalty cost}), przybierającej postać ekwiwalentu czasu doliczanego do całkowitego kosztu podróży za każdą zmianę linii. W badaniach empirycznych nad sieciami transportowymi często przyjmuje się zryczałtowaną wartość karną (np. 5 minut dodatkowego obciążenia) za sam fakt wymuszenia na pasażerze transferu~\cite{LUO2019102505}.

% \subsection{Uogólniony Koszt Podróży (GTC)}

% Agregacja omówionych wyżej składowych pozwala na wyznaczenie kompleksowej miary syntetycznej, jaką jest Uogólniony Koszt Podróży (ang. \textit{Generalized Travel Cost} – GTC). W analityce układów transportowych dopiero tak wyliczony koszt stanowi poprawną metodologicznie bazę do estymacji parametru impedancji w modelach grawitacyjnych, ponieważ obiektywnie odwzorowuje faktyczną barierę wysiłku pokonywania przestrzeni przez człowieka. Bez wdrożenia modelu całego łańcucha podróży i funkcji GTC do analiz interakcji, pomiary dostępności rynków pracy czy usług mogą prowadzić do błędnych wniosków, oderwanych od logiki funkcjonowania systemów operacyjnych i rzeczywistych zachowań przestrzennych podróżnych~\cite{LUO2019102505}.

% % Precyzyjna ocena dostępności w transporcie publicznym wymaga modelowania całego łańcucha podróży (Whole Travel Chain), a nie tylko czasu jazdy w pojeździe. Konieczne jest włączenie tzw. pierwszej i ostatniej mili (access/egress), czyli czasu dojścia pieszego od punktu początkowego do przystanku oraz od przystanku końcowego do celu.

% % Całkowity koszt podróży, wyrażany jako Uogólniony Koszt Podróży (GTC – Generalized Travel Cost), jest sumą czasu dojścia, czasu oczekiwania na przystanku (często szacowanego jako połowa interwału kursowania), czasu jazdy oraz czasu przesiadek. Co istotne, w funkcjach impedancji stosuje się penalizację przesiadek i oczekiwania, aby odzwierciedlić większą uciążliwość tych elementów dla pasażera. Na przykład, każda przesiadka może być modelowana jako dodatkowy koszt czasowy (np. 5 minut kary) dodawany do rzeczywistego czasu transferu.

% % \section{Sprawiedliwość transportowa i mapowanie społeczne}

% % Kompleksowa ocena systemów publicznych, wykraczająca poza ramy analityki topologicznej i uogólnionego kosztu podróży, wymaga osadzenia uzyskanych wyników w kontekście struktury społecznej. W nowoczesnej inżynierii transportu i planowaniu przestrzennym fundamentalne staje się rozróżnienie pomiędzy pojęciami równości (ang. equality) a sprawiedliwości lub słuszności transportowej (ang. justice / equity). Równość definiowana jest najczęściej jako czysto geometryczne, równomierne rozmieszczenie infrastruktury lub zapewnienie identycznego poziomu usług dla ogółu populacji, co w praktyce ignoruje zróżnicowanie fizycznych i ekonomicznych potrzeb mieszkańców. Sprawiedliwość transportowa poszukuje z kolei obiektywnej uczciwości w dystrybucji dostępu do mobilności, kładąc szczególny nacisk na systematyczne i bezstronne traktowanie jednostek oraz społeczności historycznie marginalizowanych i niedostatecznie obsługiwanych przez sieć. Oznacza to kategoryczne odejście od paradygmatu homogenicznego pasażera na rzecz uwzględnienia faktu, że obciążenia oraz koszty wynikające z przemieszczania się są wewnątrz tkanki miejskiej dystrybuowane wysoce asymetrycznie~\cite{KHALIFE2023160, elavating}. 

% % Weryfikacja tego zjawiska na poziomie modelowania matematycznego wymusza zastosowanie rygorystycznej segmentacji demograficznej. W podwójnie ograniczonych modelach grawitacyjnych kluczową rolę w definiowaniu funkcji impedancji odgrywa wykładniczy lub potęgowy współczynnik zaniku oporu przestrzennego. Empiryczne analizy przepływów wykazały, że wartość tego parametru nie jest uniwersalna, lecz wykazuje silną korelację ze statusem materialnym badanej grupy~\cite{capetown, gravity2}. Konieczna staje się zatem odrębna kalibracja współczynnika w zależności od stopnia zamożności. Osoby o niższych dochodach, dysponujące mocno ograniczonym budżetem rozporządzalnym, wykazują wyższą wrażliwość na koszty finansowe podróży, co zmusza je do akceptacji dłuższego czasu dojazdu w zamian za optymalizację wydatków. Z kolei grupy zamożniejsze charakteryzują się diametralnie inną krzywą – ich wrażliwość na straty czasowe jest wyższa, co modyfikuje skłonność do pokonywania oporu przestrzeni~\cite{capetown}.

% \section{Sprawiedliwość transportowa i mapowanie społeczne}

% Kompleksowa ocena publicznych systemów transportowych, wykraczająca poza ramy analityki topologicznej i wyznaczanie uogólnionego kosztu podróży, wymaga osadzenia uzyskanych wyników w kontekście struktury społecznej. W nowoczesnej inżynierii transportu i planowaniu przestrzennym istotne staje się rozróżnienie pomiędzy pojęciami równości (ang. \textit{equality}) a sprawiedliwości transportowej (ang. \textit{transport equity}). Równość definiowana jest najczęściej jako czysto geometryczne, równomierne rozmieszczenie infrastruktury lub zapewnienie identycznego poziomu usług dla ogółu populacji, co w praktyce często pomija zróżnicowanie fizycznych i ekonomicznych potrzeb mieszkańców. Sprawiedliwość transportowa poszukuje z kolei obiektywnej proporcjonalności w dystrybucji dostępu do mobilności, kładąc szczególny nacisk na systematyczne wsparcie społeczności niedostatecznie obsługiwanych przez sieć. Stanowi to wyraźne odejście od paradygmatu homogenicznego pasażera na rzecz uwzględnienia faktu, że obciążenia oraz koszty wynikające z przemieszczania się są wewnątrz tkanki miejskiej dystrybuowane asymetrycznie~\cite{KHALIFE2023160, elavating}.

% Weryfikacja tego zjawiska na poziomie modelowania matematycznego wymusza zastosowanie segmentacji demograficznej. We wprowadzonych wcześniej podwójnie ograniczonych modelach grawitacyjnych kluczową rolę odgrywa współczynnik zaniku oporu przestrzennego. Empiryczne analizy przepływów wskazują, że wartość tego parametru wykazuje silną korelację ze statusem materialnym badanej grupy~\cite{capetown, gravity2}. Z tego względu konieczna staje się jego odrębna kalibracja w zależności od profilu dochodowego. Osoby o niższym statusie materialnym, dysponujące ograniczonym budżetem rozporządzalnym, wykazują wyższą wrażliwość na koszty finansowe podróży (co wiąże się bezpośrednio z poruszaną wcześniej kwestią przystępności cenowej), zmuszając je nierzadko do akceptacji dłuższego czasu dojazdu w celu optymalizacji wydatków. Z kolei grupy zamożniejsze charakteryzują się odmienną krzywą popytu – ich wrażliwość na straty czasowe jest wyższa, co modyfikuje skłonność do pokonywania oporu przestrzeni~\cite{capetown}.

% % Współczesne badania nad dostępnością kładą nacisk na rozróżnienie między równością (*equality*) a sprawiedliwością (*justice* lub *equity*). Równość oznacza zapewnienie identycznego poziomu usług dla wszystkich, podczas gdy sprawiedliwość transportowa uwzględnia zróżnicowane potrzeby grup defaworyzowanych i dąży do wyrównywania szans.

% % W modelowaniu przekłada się to na konieczność segmentacji demograficznej. Tradycyjne modele grawitacyjne (Hansena) często ignorują różnice społeczno-ekonomiczne. Wprowadzenie segmentacji pozwala na kalibrację współczynnika zaniku $\beta$ w zależności od grupy dochodowej – osoby o niższych dochodach mogą być bardziej wrażliwe na koszty pieniężne, a mniej na czas, lub odwrotnie, w zależności od dostępności alternatyw. Ponadto, dla osób z niepełnosprawnościami (np. poruszających się na wózkach inwalidzkich), bariery pozamaterialne, takie jak brak obniżonych krawężników czy wind, drastycznie redukują dostępność sieci. Badania wykazują, że dostępność do miejsc pracy dla osób na wózkach może stanowić jedynie ułamek dostępności oferowanej ogółowi populacji, co jest kluczowym elementem wykluczenia transportowego.


\chapter{Wykorzystane zbiory danych i ich opracowanie}\label{chapter:dane}

\printbibliography

\end{document}

% %%% Local Variables:
% %%% mode: latex
% %%% TeX-master: t
% %%% coding: latin-2
% %%% End: